\documentclass[10pt,a4paper]{amsart}
\usepackage[margin=19mm]{geometry}
\usepackage{amsmath,amssymb,mathtools}
\usepackage[scr=rsfs,cal=euler]{mathalfa}
\usepackage{xfrac}
\usepackage[unicode,hidelinks,bookmarks=false]{hyperref}
\newcommand{\ie}{i.e.\ }
\newcommand{\cf}{cf.\ }
\newcommand{\resp}{respectively\ }
\newcommand{\Subsection}{\subsection}
\providecommand{\nobreakdash}{\nobreak\hbox{-}}
\setlength{\emergencystretch}{2em}
\multlinegap0pt
\usepackage{mathtools}
\newtheorem{lemma}{Lemma}
\newcommand{\ti}{\tilde}
\title{Automorphism groups and Homogeneous fibrations on projective varieties with nef anticanonical divisors}
\author{Zhan Li, Jinsong Xu}
\date{Source: arXiv:2609.09976v1 (CC BY 4.0)}
\begin{document}
\maketitle

\noindent\emph{Source and licence.} Automorphism groups and Homogeneous fibrations on projective varieties with nef anticanonical divisors. Version arXiv:2609.09976v1 (CC BY 4.0), distributed by arXiv under the Creative Commons Attribution 4.0 International licence, http://creativecommons.org/licenses/by/4.0/, as recorded in the arXiv licence metadata for this exact version and shown on the abstract page https://arxiv.org/abs/2609.09976v1. Full licence notice: \texttt{licenses/arxiv-2609.09976-CC-BY-4.0.txt}.\par\smallskip

\noindent\textit{Editorial scope.} Counted unit: the article's lemma field (source0.tex lines 301--306). The article's complete original proof of it is delivered with it (source0.tex lines 307--328, one \texttt{proof} environment of 22 lines). No mathematics is written, repaired or re-derived: the statement and the proof are the article's own lines, in the article's own notation, and no retained line is substituted or re-flowed.  No further source-local context is needed: every cross-reference of every retained line resolves inside the two delivered spans.  Nothing was omitted from any retained span, so the package carries no omission record.  No retained line contains a citation, so the wrapper carries no bibliography and no bibliography entry.  No retained line contains a figure environment, an image-inclusion command or drawing code, so no image is delivered and the package declares no asset dependency.  Editorial formatting only, in addition: an AMS \texttt{amsart} preamble that restates the article's own declarations and macros used by the retained lines, namely the environments \texttt{lemma} on separate counters, together with the standard packages those lines need.  The retained environments print in this document as Lemma 1 in order of appearance, whereas the article prints the same items with the numbering of its own full document.  The complete native source of the article as distributed in the arXiv e-print for this exact version is bundled unchanged as \texttt{sources/209864/source0.tex}.
\begin{lemma}\label{lem: lift vector field}
    If $f: X \coloneqq \ti S \times^\rho F \to S$ is a locally constant fibration, then the natural map 
    \[
    T_f: H^0(X, T_X) \to H^0(X,f^*T_S)=H^0(S, T_S)
    \] is surjective. 
\end{lemma}

\begin{proof}
    Choose $v \in H^0(S, T_S)$, then $v$ naturally lifts to a vector field $\ti v \in H^0(\ti S, T_{\ti S})$ such that
    \[
    \ti v(\gamma \cdot \ti t) = v(x) \in T_{\ti S, \gamma\ti t} \simeq T_{S, t}
    \] for any $\gamma\in \pi_1(S)$, $\ti t \in \ti S$ and $t = \pi(\ti t)$. Then, $\ti v$ trivially lifts to a vector field $\ti w$ on $\ti S \times F$ such that
    \[
    \ti w(\ti t, y)= \ti v(\ti t) \times \{0\} \in T_{\ti S, \ti t} \times T_{F,y},
    \] where $\ti t\in\ti S$ and $y\in F$.

    We claim that $\ti w$ is $\pi_1(S)$-invariant. That is,
    \[
    \ti w(\gamma \cdot \ti t, \rho(\gamma)(y)) = \ti w(\ti t, y)
    \] for any $\gamma\in \pi_1(S)$. This holds as 
    \[
    \ti w(\gamma\cdot \ti t, \rho(\gamma)(y)) = \ti v(\gamma\cdot \ti t) \times \{0\}= v(t) \times \{0\}
    \] and
    \[
    \ti w(\ti t, y) = \ti v(\ti t) \times \{0\}= v(t) \times \{0\}
    \] by construction.

    Therefore, $\ti w$ descents to a vector field $w$ on $X$. It is straightforward to see that $T_f(w)=v$.
\end{proof}

\end{document}
